\section{Discussion}
\label{sec:discussion}

\subsection{What an additive all-atom energy makes available}
\label{sec:whatchanges}

The sampling half of natural-move Monte Carlo is unchanged by any of this. What
changes is what an accepted structure is worth. A knowledge-based potential
returns a ranking: it says which of two models is preferred, on a scale fitted
to that comparison and to the representation it was fitted on
\cite{minary2008foldspace,leaverfay2011rosetta}. For structure prediction, and
for refinement where an experimental restraint carries most of the signal
\cite{zhang2012multiscale,chang2020rnacryoem}, that is sufficient. It is not
sufficient for anything that needs the energy itself. Two such numbers computed
for two different molecules do not lie on the same axis, nothing computed from
the energy on an absolute scale is available, and a structure cannot pass
between the Monte Carlo engine and a molecular dynamics run without the sampled
quantity changing definition on the way.

With eight parameter families evaluated to the tolerances of
Sections~\ref{sec:matrix} and~\ref{sec:energies}, those obstructions are lifted
for the chemistry and the geometries tested, and for nothing beyond them. The
claim the data supports is narrow and should be stated narrowly: handed a
published deck, the engine reproduces that deck's energy and its gradient to
the tolerance reported, whichever deck it is. The four AMBER nucleic decks are
substantially different parameter sets, and the engine follows each of them
without being retuned for any. What that does not license is a statement about
how far apart the decks put a molecule. A Fourier torsion sum, $\sum
(V_n/2)[1+\cos(n\phi-\gamma)]$, does not vanish at any reference state, and a
refit moves its additive constant as freely as it moves its shape; the
quantity that would carry that argument is $\Delta\Delta E$ between at least
two conformations of the same molecule under two decks, and it is not computed
here.

What the engine can now be asked, for the systems tested, is a different class
of question. It can be run at a stated temperature against a named force field,
so that an ensemble it produces and an ensemble a molecular dynamics code
produces under the same parameters are ensembles of the same function; and a
configuration produced by a collective move can be continued under molecular
dynamics, or taken from one, without the energy being redefined in between.
Neither is demonstrated in this paper. What is demonstrated is the
precondition: the energy and its gradient are the published ones, at the
geometries and on the chemistry tested.

\subsection{What the extension to eleven systems adds, and what it does not}
\label{sec:elevenadds}

The two reference systems of Section~\ref{sec:energies} settle whether the
engine evaluates a deck correctly, term by term, against three programs. The
other nine settle a different question: whether it does so on the chemistry a
user will actually bring. A deck is a few hundred to a few thousand parameter
rows and a 56-atom peptide reaches a fraction of them. Ubiquitin, BPTI and the
MHC heavy chain between them select every amino-acid type ff14SB and ff19SB
carry, with charged, aromatic, histidine and proline residues, disulfides and
charged termini; the dodecamer and the tetraloop select all eight bases and both
$5'$-OH and $3'$-OH ends; the two complexes select a bound peptide in a groove
and a protein--RNA interface. That every one of those cells lands within a few
thousandths of a \kcal{} of the references, and that the two that land further
away resolve to the atom order of a single improper, is what the extension
adds.

It also sharpens what the residuals mean. On the two reference systems the
angle term carried the largest residual and it was attributed to the reference
topology (F-12). Seven further systems now carry the same residual with the
same sign, growing with the number of angles, which is what a rounding of
$\pi/180$ in the topology builder predicts and what a defect in the engine's
bend routine would not. The improper cells make the complementary point: the
one large disagreement in the matrix is not a parameter or an expression but a
convention, the order in which four atoms are listed, and it appears where the
geometry departs from planarity and nowhere else. Cross-engine agreement
certifies the evaluation path. It does not certify the path that produced the
parameters, and both of the residuals that survive at the top of the matrix
live there.

What the extension does not add is a per-term comparison against all three
references for every cell. The totals are against two or three programs on 20
of the 26 cells and against OpenMM alone on the other six; the per-term
resolution in Figure~\ref{fig:energy-agreement}(b) exists only where the output
files carry it, and eleven OpenMM comparisons print a single non-bonded sum.
Nor does it add a second conformation per system. The extension widens the
chemistry and holds the geometry fixed, and the failure mode F-21 demonstrates
remains open on every cell.

\subsection{What changed in the engine, and what did not}
\label{sec:baseline}

A reader is entitled to know how much of this is new code, because that is
where the risk of a new defect lives. The honest answer is that most of it is
not. Both CHARMM-specific terms were carried on machinery the engine already
had. The 1497 Urey--Bradley entries in CHARMM36 reduce to 810 unordered
triples, of which 807 are written into the bond slot of the existing bend
record, which \texttt{bend.cpp} already evaluates as
$\tfrac{1}{2}k(r_{13}-r_0)^2$, the functional form CHARMM asks for; three
triples are dropped and counted, because four P--ON2--P variants carry negative
force constants that CHARMM permits and the deck reader rejects, and they are
pyrophosphate linkages that do not occur in a phosphodiester backbone (F-16).
The 188 harmonic impropers are carried on the existing torsion record with
\verb|\pot_type{harm}|, which the deck reader expands onto the same power series
it uses for the cosine form. That expansion is a genuine approximation, exact
to four figures at $10^{\circ}$ from equilibrium and $-7.7$\,\% at
$60^{\circ}$, and the ten-digit agreement reported for the improper term is at
a total of $1.0\times10^{-3}$~\kcal{}, where the series and the quadratic
coincide by construction.

Where the defects were was the converter and, in a few places, the engine. The
CHARMM XML splits the non-bonded term between two elements, charge in
\texttt{NonbondedForce} and $\sigma$, $\epsilon$ in
\texttt{LennardJonesForce}; the converter read only the first, so an entire
CHARMM deck carried $\epsilon = 0$ on all 85\,078 pair entries. The
correction-map analytic force was identically zero at $\phi = \pi$, because a
$\sin\phi$ clamp replaced a vanishing denominator and left a vanishing
numerator (F-21); it is now written in the form with no $\sin\phi$ in it. Forty
arrays allocated with \verb|new[]| were grown with \verb|realloc| (F-19), which
ran on every committed job and was invisible because the relevant sanitizer
check is off by default on macOS. And the deck reader now carries the
per-molecule torsion label and chains the residues it builds; the released
3.9.1 refuses the decks that need the first and returns energies tens of
thousands of \kcal{} from the references without the second.

The novelty is not all-atom Monte Carlo on an additive force field. That
exists: CAMPARI with ABSINTH, MCPRO and BOSS, ProtoMS and Cassandra all evaluate
all-atom additive potentials under Metropolis Monte Carlo, and concerted-rotation
move sets have been run in OPLS for both protein and nucleic acid backbones
\cite{ulmschneider2003concerted,ulmschneider2004concertedna,vitalis2009mcmethods}.
What is new here is the pairing of that energy with the nested natural-move
hierarchy and its stochastic chain closure
\cite{minary2010chainclosure,sim2012naturalmoves}, so that one proposal can act
at the scale of a domain and the next at the scale of a sugar pucker while the
accepted energy stays a published one, and the delivery of that pairing to a
user who has never seen the engine's input syntax. The claim is the combination,
and it should be read that way.

\subsection{Three references, and what the third one pays for}
\label{sec:threerefs}

Agreement with a single reference program is a weaker claim than it looks. It
establishes that two codes treat a parameter file the same way, which confounds
correct evaluation with shared convention: the same choice of 1--4 scaling, the
same grouping of impropers, the same charge-to-energy constant. Three
independently written programs break that confound in both directions. Where
they agree with one another, the target is not in dispute and a residual
against any of them is a measurement of the implementation. Where they
disagree, the disagreement is itself a number, and it bounds what any
cross-engine claim can resolve. On 1efs the bound is $2.6\times10^{-5}$~\kcal{}
on Lennard-Jones and $3.4\times10^{-5}$ on electrostatics, which is what OpenMM
and GROMACS return against each other when handed the same parameters and the
same coordinates through the ordinary conversion route.

Each of the two convention differences reported here was invisible to some
subset of the three. The GROMACS route is the only one that could not carry
ff19SB's per-residue correction maps, and running it put a number on what that
costs: 0.30~\kcal{} on a four-residue peptide and 8.4~\kcal{} on ubiquitin, from
applying one residue's map to every residue. sander, separately, is the only one
of the three that does not share the modern charge-to-energy constant, and a
comparison against OpenMM and GROMACS alone would have shown MOSAICS's own
$8.2\times10^{-7}$ offset while leaving the $3.5\times10^{-5}$ carried by a
widely used engine undetected. Neither is a defect in the engine that exposed
it, and neither would have been found by adding a fourth code from the same
family. The independence is in the code and not in the input, and that limit is
worth stating because it is where the argument is weakest: all three references
read a topology written by one builder, so an error in the builder appears
identically in all three, and F-12 is exactly that case. The three references
are also not of equal weight: six combinations against three engines, fourteen
more against two, and six against OpenMM alone are not the same strength of
claim, and Table~\ref{tab:matrix} keeps them apart.

\subsection{The charge-to-energy constant is a general hazard}
\label{sec:constantgeneral}

The constant that converts $q_iq_j/r$ into \kcal{} is not the same number in
every program, and any group differencing energies across programs is exposed
to that. sander uses $18.2223^2 = 332.052217$, an older value carried in the
way it stores topology charges; OpenMM and GROMACS both return $332.063713$
from the same two-charge measurement. The relative gap is $3.5\times10^{-5}$,
which on the 1efs electrostatic term of 1857.7~\kcal{} is 0.064~\kcal{}, roughly
360 times the largest residual this work is trying to resolve and about a tenth
of $kT$ at 300~K. The constant enters multiplicatively, so it does not cancel
when two conformations of the same molecule are compared: it biases the
electrostatic part of that difference by the same $3.5\times10^{-5}$, and its
absolute size grows with the electrostatic energy, which is to say with system
size and with net charge. It also has nothing to do with the AMBER family: the
same offset accounts for the CHARMM36 residuals on every one of the five
CHARMM36 cells, and once it is removed those cells are the closest agreement in
this work, at a few parts in $10^{12}$ of the term. The practical
recommendation is to measure the constant instead of taking it from
documentation, which is a few minutes of work in any engine that will read a
topology, and to state which constant a comparison was corrected to, because
the choice is worth $1.7\times10^{-6}$~\kcal{} on this system.

\subsection{What these numbers do not establish}
\label{sec:limits}

The first limit is the one conformation. Every run sets all proposal widths to
zero, so a term that happens to vanish at that geometry is untested. The capped
peptide sits at the minimum of its out-of-plane terms and both MOSAICS and
GROMACS return zero for them (F-11). More seriously, the correction-map
analytic force was identically zero, 0 of 168 components, because a peptide
built extended sits at $\phi = \psi = \pi$, the singular point of the
implemented expression, and \emph{every energy comparison on that system passed
while that was true} (F-21). It took a gradient sweep to find it. The same
failure mode is available to any term whose parameter is wrong but inactive, or
whose phase is wrong but degenerate, at the one geometry tested, on any of the
26 cells. Repeating these single points over 50 to 100 decorrelated snapshots
per system, as Shirts \emph{et al.}\cite{shirts2017engines} do, would close it,
and it is cheap on vacuum systems of this size. It has not been done here.

The second is that the force sweeps cover four systems and the term-by-term
comparison two. The nine added systems carry totals, and per-term differences
where the output files resolve them, but no force sweep and no staged sweep.
The correction map is exercised at the force level on one system, because the
ff19SB peptide is the only committed force input that loads a map database and
no nucleic parameter set here carries one (F-22).

One committed structure is geometrically defective and the defect is not
repaired. In \path{examples/pure_terminal_controls/dna_acgt.pdb}, residue DG~3,
H5$'$ and H5$''$ each sit 1.089 and 1.090~\AA{} from C5$'$, which is right, and
0.060~\AA{} from each other, which makes the H--C--H angle $3.15^{\circ}$
instead of $109.5^{\circ}$ (F-17). The offending pair is 1--3, which the
non-bonded sum excludes, so no energy check sees it as an error; the symptom is
that this 126-atom 4-mer returns a bend energy of 508.62~\kcal{} under the AMBER
decks against 220.56 for the 818-atom heteroduplex, and 412.79 under CHARMM36
against 166.36 for the RNA 4-mer of the same size. The cross-engine comparison
is unaffected, because every program reads the same coordinates and they
reproduce each other on them. The file was pinned in a test and left alone
because committed reference tables were computed from it; repairing it and
regenerating those tables is the right fix and has not been done.

Three scope limits belong with these. CHARMM36 for protein is not available in
MOSAICS, and the obstruction is a file format and not a physical one: its
correction maps are selected partly by the residue that follows the one being
corrected, while the engine's map-selection record keys on the residue itself
(F-09), and no claim about the CHARMM36 protein parameters
\cite{best2012charmm36,huang2017charmm36m} appears anywhere in this work. Two
decks cannot be regenerated from their sources (F-20), so the seventh decimal
of the OL15/OL3 and OL21/OL3 bonded terms is reporting transcription; and the
OL24 parameter files come from a package whose redistribution terms are its
authors', so a reader who does not hold that package cannot rebuild those three
cells from a clone. And parmbsc1 has no terminal deck and parmbsc0 does not run
on a free-ended duplex, which is why the dodecamer has no cell under either.

There is a gap between the code path validated here and the code path a
production run executes, and it is narrower than that sentence usually implies.
MOSAICS scores a proposal by recomputing the whole potential, so the term
routines checked here are the term routines a proposal is judged by. What no
single point exercises is a choice of pair set: \texttt{inter\_list} is
\texttt{none} in every committed input, so the link-list and neighbour-list
branches never ran. Closing that costs one single point per branch with a list
cutoff longer than the molecule, differenced against the numbers reported here
to the summation-order tolerance of $5.5\times10^{-11}$~\kcal{}. Neither run is
in this paper.

Single-point energies and finite-difference force checks fix the value of the
energy function at a geometry and its gradient there. They say nothing about
whether the move set, the stochastic chain closure and the acceptance rule
together generate the Boltzmann distribution of that energy
\cite{minary2010chainclosure,hatch2025mcbestpractices}; that is a separate
validation with separate diagnostics, and it is not attempted here. The cost
figure of Section~\ref{sec:cost} is likewise not an efficiency result: a
per-step time on one machine compares nothing with anything, and the quantity
that would, sampling per unit cost against a molecular dynamics baseline on the
same system and the same force field, is the next measurement and not this
one. Nor do these numbers speak to whether ff19SB or OL21/OL3 is right about a
molecule, which is measured against experiment and not against another program
\cite{lindorfflarsen2012validation,beauchamp2012benchmark,love2023amberdna}.
And everything was run in vacuum at dielectric~1 with no cutoff, so no cutoff
scheme, switching function or implicit-solvent term is covered by any number in
this paper.

\subsection{What the Studio does and does not establish}
\label{sec:studio-limits}

The Studio's verification is of the engine it carries, on one documented run.
The cross-check shows that the WebAssembly build follows the native one through
2,040 minimisation steps to the same energies at the precision the energy file
carries and to 0.0045~\AA{} in the final structure; it does not show that every
protocol behaves identically in the browser, and two protocols are known to
differ by construction, because the browser engine is single-threaded and
parallel-tempering replicas take turns. The record binds what ran to what was
produced; it does not make the run correct, and the validation gate on the
Prepare step checks the ordinary energy terms of the residue types present
against a committed reference, which is the baseline of this paper and not
more. No external user has run the software at the time of writing, and
proteins are prepared from coordinates rather than built.

\subsection{Next}
\label{sec:next}

The intended application is conformational sampling of peptide--MHC complexes,
and it is the next piece of work rather than a result reported here. The system
suits the move set: a bound peptide with a small number of internal degrees of
freedom in a groove, flanked by domains that move as units, which is the
granularity natural moves were built for \cite{sim2012naturalmoves}, and
hierarchical natural-move Monte Carlo has already been applied to peptide
dissociation from MHC class~I \cite{knapp2016pmhc}. What exists in this work is
the implementation side of that study: the complex evaluates under both
protein force fields to within 0.005~\kcal{} of two references at 6,139 atoms,
the staging transform was checked on a peptide--MHC fixture at the gradient
level, and the Studio will prepare and run it. No sampling run of any kind on
that system appears in the evidence.

Several things should be settled before such a study is worth reporting, and
each is small and named. The single points should be repeated over an ensemble
of decorrelated configurations per system, the only cheap defence against the
failure mode F-21 demonstrates. The force sweeps should be extended to the nine
added systems. The two large cells that are not resolved into terms should be.
The link-list and neighbour-list branches of the non-bonded sum need the two
single points described above. The defective control structure should be
repaired and its reference tables regenerated (F-17); the two decks that cannot
be regenerated should be rebuilt (F-20); the peptide needs a deliberately
distorted copy so that the out-of-plane terms and the CHARMM improper series
are exercised away from their minima (F-11). Beyond the fixtures, the certified
regime has to be widened from vacuum with no cutoff to whatever solvent
treatment a production run will use, which is a further round of the same
term-by-term comparison and not an extrapolation from these numbers. The
criterion for judging the study that follows should not be whether the Monte
Carlo engine reproduces a molecular dynamics trajectory. It should be whether
the engine reaches configurations that molecular dynamics does not reach in
comparable time, at energies that mean the same thing in both. The present
paper speaks only to the second half of that sentence, and only on the
chemistry and the geometries it tested.
